% =============================================================================
%  Phase 3: EDA & Feature Engineering
%  Rubric: docs/spec/phase3_details.pdf (no page limit; "concise").
%  Source of every number: notebooks/feature_engineering.ipynb (executed outputs).
%  Figures: reports/figures/p3_*.png, written by that notebook.
%  Lines marked "% CHECK:" are review notes for the authors, not report text.
%  Build: latexmk -pdf -outdir=build main.tex
% =============================================================================
\documentclass[11pt,letterpaper]{article}
\input{../_template/preamble}
\addbibresource{references.bib}

% ---- Metadata ---------------------------------------------------------------
\projecttitle[Bill Forecasting]{Forecasting Household Month End Electricity Bills from Smart Meter Data}
\phaselabel{Phase 3 \textperiodcentered{} EDA \& Feature Engineering}
\authors{Akinbobola Akin (3784664), Adejuwon Oluwasegun Isaac (3806766), Edehi Andrew Igue (3778767)}
\course{TME6014 -- Data Analytics}
\reportdate{November 10, 2026}
% CHECK: set to the actual submission date.

\begin{document}
\maketitleblock

% =============================================================================
\section{Introduction and Scope}
\label{sec:intro}

Predicting a household's electricity bill before the billing month ends lets the household react while
there is still time to change its usage. This phase explores the cleaned smart meter data from Phase~2 and
builds the modelling dataset for Phase~4. It returns to the research question set in Phase~1:

\begin{quote}
\textbf{On day $d$ of a billing month, how accurately can a household's month-end bill be projected, with a
prediction interval, from its smart meter readings?}
\end{quote}

On day $d$ the kWh used on days $1..d$ is already known. The model only has to predict the kWh for days
$d{+}1$ to the end of the month and add it to the known part. Three candidate values, $d \in \{7, 14, 21\}$,
are carried forward; Phase~4 chooses the final one.

\paragraph{Bill proxy} The bill is a \emph{computed proxy}, not a real invoice: kWh $\times$ \pounds 0.14228
per kWh, the published standard tariff~\cite{datastore_smartmeters}. Because the bill is kWh times a fixed
rate, every correlation ``with the bill'' in this report is identical to the correlation with kWh.

\paragraph{Scope} The analysis covers \textbf{standard-tariff (Std) households only}. Time-of-use households
paid different prices by time of day during 2013, so the flat-rate proxy misstates their bills; they are kept
for the stretch goal. The filter is applied before any statistic is computed. No household changes tariff
during the trial, so whole histories are kept or dropped (Table~\ref{tab:scope}).

\begin{table}[!ht]
  \centering
  \caption{Data used in this phase (Phase~2 outputs~\cite{kaggle_smartmeters}, December 2011 to January 2014).}
  \label{tab:scope}
  \small
  % Y{f} widths sum to 1 (one Y column).
  \begin{tabularx}{\linewidth}{Y{1} r r r}
    \tableheadrow
    \colhead{Level} & \colhead{All tariffs} & \colhead{Std only} & \colhead{Kept} \\
    \tableheadend
    Households & 5,548 & 4,428 & 79.8\% \\
    Household-months (\code{monthly.parquet}) & 106,225 & 84,807 & 79.8\% \\
    Household-days (\code{daily.parquet}) & 3,239,805 & 2,586,554 & 79.8\% \\
    \bottomrule
  \end{tabularx}
\end{table}

Every number and figure in this report comes from the notebook \code{feature\_\allowbreak engineering.ipynb}. It reads
the Phase~2 Parquet files with DuckDB and saves the modelling dataset as \code{features\_\allowbreak day\_d.parquet}
in \code{data/processed/}.

% =============================================================================
\section{Exploratory Data Analysis}
\label{sec:eda}

\subsection{Descriptive Statistics}
\label{sec:stats}

Table~\ref{tab:stats} summarises the monthly and daily variables. No value is missing. Households have a
median of 20 months of data (range 1 to 26), and 4,244 of the 4,428 have at least 12 months.

\begin{table}[!ht]
  \centering
  \caption{Summary statistics, Std households. \code{filled\_slots} counts half hours that Phase~2 filled in,
  out of about 1,460 per month.}
  \label{tab:stats}
  \small
  \begin{tabularx}{\linewidth}{X r r r r r}
    \tableheadrow
    \colhead{Variable} & \colhead{Mean} & \colhead{Median} & \colhead{SD} & \colhead{Min} & \colhead{Max} \\
    \tableheadend
    Monthly kWh & 315.96 & 250.29 & 261.81 & 1.29 & 6,385.89 \\
    Monthly bill (\pounds) & 44.95 & 35.61 & 37.25 & 0.18 & 908.58 \\
    Monthly mean temperature (\textdegree C) & 11.01 & 10.96 & 5.04 & 3.80 & 20.04 \\
    Filled half hours per month & 2.30 & 0 & 12.83 & 0 & 146 \\
    Daily kWh & 10.36 & 7.96 & 9.35 & 0.001 & 332.56 \\
    \bottomrule
  \end{tabularx}
\end{table}

Consumption is strongly right-skewed (Figure~\ref{fig:hist}): the mean is well above the median and the largest
month is about 25 times the median. Bills have exactly the same shape, since they are kWh times a constant.

\begin{figure}[!ht]
  \centering
  \includegraphics[width=0.75\linewidth]{p3_monthly_kwh_hist.png}
  \caption{Monthly kWh per household-month. The axis stops at the 99.5th percentile so the long tail does not
  flatten the plot.}
  \label{fig:hist}
\end{figure}

\subsection{Seasonality and Temperature}
\label{sec:season}

Consumption follows a clear yearly cycle that mirrors temperature (Figure~\ref{fig:trend}). Mean winter
consumption is 379.49~kWh per household-month against 256.22~kWh in summer, 48.1\% higher (spring 323.74,
autumn 308.92). Across the 26 calendar-month averages, temperature and kWh correlate at $r = -0.925$, but at
the household-month level the correlation is only $-0.186$. Differences \emph{between} households (size,
heating type, occupancy) are much larger than the seasonal swing \emph{within} a household. Temperature
explains when a household uses more, not which household uses more.

Two details matter for modelling. The first months of the trial hold few households (61 in December 2011),
so their averages are noisy; from June 2012 every month has more than 2,800. And February 2013 dips below
its neighbours even though it was colder than January, because February has only 28 days: monthly totals
partly measure month length, which the day-$d$ features handle by working with daily means and day counts.

\begin{figure}[!ht]
  \centering
  \includegraphics[width=0.82\linewidth]{p3_monthly_trend_temp.png}
  \caption{Mean kWh per household (top) and mean temperature (bottom) for each calendar month. Two panels
  rather than two y-axes, so neither scale exaggerates the relationship.}
  \label{fig:trend}
\end{figure}

\subsection{Household Groups and History}
\label{sec:groups}

ACORN classifies neighbourhoods by income and lifestyle. Affluent households use the most and Adversity
households the least (Table~\ref{tab:acorn}), but the distributions overlap heavily
(Figure~\ref{fig:acorn}): the group shifts the typical level and leaves most of the variation unexplained.
The Unknown group has only 37 households, so its high mean is unstable (its median is below Affluent's). These
differences are descriptive and do not establish causal relationships.

\begin{table}[!ht]
  \centering
  \caption{Monthly kWh by ACORN group, Std households.}
  \label{tab:acorn}
  \small
  \begin{tabularx}{\linewidth}{X r r r r}
    \tableheadrow
    \colhead{ACORN group} & \colhead{Households} & \colhead{Household-months} & \colhead{Mean kWh} & \colhead{Median kWh} \\
    \tableheadend
    Affluent & 1,700 & 33,068 & 361.69 & 266.70 \\
    Comfortable & 1,179 & 22,381 & 313.42 & 267.34 \\
    Adversity & 1,512 & 28,692 & 263.86 & 224.09 \\
    Unknown & 37 & 666 & 375.49 & 238.00 \\
    \bottomrule
  \end{tabularx}
\end{table}

\begin{figure}[!ht]
  \centering
  \includegraphics[width=0.82\linewidth]{p3_acorn_box.png}
  \caption{Monthly kWh by ACORN group (box: middle 50\%; line: median; whiskers: 1.5 IQR; outliers hidden).}
  \label{fig:acorn}
\end{figure}

A household's months are not always consecutive, because Phase~2 dropped incomplete months. Of 80,379 pairs
of consecutive records, 240 (in 175 households) skip at least one month. Lag features must check that the
earlier record is the exact calendar month required, or they silently use the wrong month.

\subsection{Outliers}
\label{sec:outliers}

The first pass flagged months above a pooled Tukey fence ($Q_3 + 1.5\,\mathrm{IQR} = 736.19$~kWh). On Std
households this flags 4,657 household-months (5.49\%) in 676 households. But 58.3\% of those flags belong to
households whose \emph{median} month is already above the fence. The pooled check finds large households,
not unusual months.

To find unusual months we compare each month with the household's own history, using the robust z-score
$0.6745\,(x - \tilde{x}_h)/\mathrm{MAD}_h$, where $\tilde{x}_h$ is the household's median monthly kWh and
$\mathrm{MAD}_h$ its median absolute deviation. Median and MAD are not distorted by the outliers being sought.
Months with $|z| > 3.5$ are flagged, for the 4,393 households with at least 6 months. Table~\ref{tab:outliers}
compares the two checks.

\begin{table}[!ht]
  \centering
  \caption{Pooled versus within-household outlier checks.}
  \label{tab:outliers}
  \small
  % Y{f} widths sum to 1 (one Y column).
  \begin{tabularx}{\linewidth}{Y{1} r r}
    \tableheadrow
    \colhead{Measure} & \colhead{Pooled IQR} & \colhead{Within household} \\
    \tableheadend
    Household-months flagged & 4,657 (5.49\%) & 1,731 (2.04\%) \\
    Households with a flag & 676 & 932 \\
    Flagged high / low & 4,657 / 0 & 1,303 / 428 \\
    Months flagged by both checks & \multicolumn{2}{c}{274} \\
    \bottomrule
  \end{tabularx}
\end{table}

The two checks find mostly different months. Within-household high flags cluster in December (348) and
January (324), consistent with a household switching on electric heating; low flags peak in August (78),
consistent with a home left empty during the summer holidays. These are real readings, so they are kept:
high consumption does not by itself indicate an error, and an alert system has to cope with exactly these
months. They are also the months a forecast will miss most, which is why the forecast needs a prediction
interval. The within-household check uses each household's full history, including later months, so it is a
data check only and never a model input.

\subsection{Daily Patterns}
\label{sec:daily}

The day-$d$ task is built from daily data, which the first pass did not use. To compare days across
households, each day's kWh is divided by that household's mean day in the same month. Weekends run 3.5\%
above that mean and normal weekdays 1.5\% below; Sunday is the highest day (+5.7\%) and bank holidays sit in
between (+1.8\%). Daily kWh falls steadily as temperature rises and levels off above about 17\textdegree C,
when there is no heating left to switch off (Figure~\ref{fig:daily}). The weekend and holiday effects are
small but consistent, so the number of such days still to come in the month carries a little information about
the remaining kWh.

\begin{figure}[!ht]
  \centering
  \includegraphics[width=\linewidth]{p3_daily_patterns.png}
  \caption{Left: daily kWh per household-day (to the 99.5th percentile). Centre: kWh by day of week relative to
  the household-month mean (bank holidays excluded). Right: mean daily kWh by daily mean temperature.}
  \label{fig:daily}
\end{figure}

\subsection{Month-to-Date Information and Benchmarks}
\label{sec:benchmarks}

Figure~\ref{fig:share} shows the share of the month's kWh already used by day $d$. The median follows the
``every day equal'' line ($d$ divided by the days in the month) almost exactly: 22.8\%, 45.7\% and 68.6\% at
$d = 7, 14, 21$, against 23.0\%, 45.9\% and 68.9\%. Scaling the month-to-date total up to a full month is
therefore a sensible baseline. The spread is the difficulty: at $d = 14$ one household-month in ten has used
less than 40.7\% of its total and one in ten more than 51.4\%. That spread is what the model must explain
from history, calendar and weather.

\begin{figure}[!ht]
  \centering
  \includegraphics[width=0.82\linewidth]{p3_share_by_day.png}
  \caption{Share of the month's kWh used by day $d$: median household-month, 10th to 90th percentile band, and
  the share if every day were equal. Dotted lines mark the candidate values of $d$.}
  \label{fig:share}
\end{figure}

This gives the first benchmark from Phase~1, the \textbf{run-rate projection}
(\code{kwh\_to\_date}~$/\,d~\times$~days in month). We compare it with the other two, \textbf{last month's bill}
and \textbf{the same month last year}, using the median absolute percentage error of the projected month
total and the mean absolute error in pounds. The median is used because a few near-empty months make
percentage errors explode. Same-month-last-year values exist only from December 2012 (most households joined
between April and July 2012), so all three are compared on the 32,491 household-months that have both a
previous month and a same month last year (Table~\ref{tab:benchmarks}, Figure~\ref{fig:runrate}).

\begin{table}[!ht]
  \centering
  \caption{Benchmark errors on the same 32,491 household-months.}
  \label{tab:benchmarks}
  \small
  \begin{tabularx}{\linewidth}{X r r}
    \tableheadrow
    \colhead{Benchmark} & \colhead{Median abs.\ \% error} & \colhead{Mean abs.\ error (\pounds)} \\
    \tableheadend
    Run-rate projection, $d = 7$ & 6.8\% & 5.04 \\
    Run-rate projection, $d = 14$ & 4.2\% & 3.23 \\
    Run-rate projection, $d = 21$ & 2.5\% & 1.93 \\
    Last month's bill & 9.3\% & 6.68 \\
    Same month last year & 12.0\% & 8.46 \\
    \bottomrule
  \end{tabularx}
\end{table}

\begin{figure}[!ht]
  \centering
  \includegraphics[width=0.82\linewidth]{p3_runrate_error.png}
  \caption{Median absolute percentage error of the projected month total by forecast day. The history-only
  benchmarks do not use days $1..d$, so they are flat.}
  \label{fig:runrate}
\end{figure}

The run-rate projection beats last month's bill from $d = 4$ and is less than half its error by $d = 14$.
This is the bar a Phase~4 model has to clear: a model that cannot beat a straight-line projection of the month
so far adds nothing. It also shows that the month-to-date total carries most of the signal, and that history
and calendar features can add most at small $d$, where the run-rate projection is weakest.

% =============================================================================
\section{Key Insights and Implications for Phase 4}
\label{sec:insights}

Table~\ref{tab:insights} lists the main findings and the modelling decision each one leads to.

\begin{table}[!ht]
  \centering
  \caption{Insights and how they guide the modelling phase.}
  \label{tab:insights}
  \small
  % Y{f} widths sum to 2 (two Y columns).
  \begin{tabularx}{\linewidth}{Y{1} Y{1}}
    \tableheadrow
    \colhead{Insight} & \colhead{Implication for Phase 4} \\
    \tableheadend
    The run-rate projection has a median error of 6.8\%, 4.2\% and 2.5\% at $d = 7, 14, 21$, well below last
    month (9.3\%) and same month last year (12.0\%). &
    Score every model against the run-rate projection at each $d$; a model is only useful if it beats it. \\
    Month-to-date usage dominates; history features correlate almost as strongly with the remaining kWh at
    $d = 7$ but fall behind by $d = 21$. &
    Predict the \emph{remaining} kWh (days $d{+}1$..end) so the known part is never re-predicted; expect history
    features to matter most at small $d$. \\
    Pooled correlation of consecutive months is 0.928, but the median within-household correlation is 0.577
    (4,366 households). &
    Much of the pooled value is ``big households stay big''. Compare pooled and per-household models (the
    secondary question) and report errors per household. \\
    Winter consumption is 48.1\% above summer; calendar-month temperature correlates with kWh at $-0.925$. &
    Keep a cyclical month encoding and temperature to date; expect them to overlap ($r = -0.89$). \\
    Usage is right-skewed and 2.04\% of months are unusual for their own household. &
    Use robust error measures (median, pounds), try a log target or robust scaling, and size the prediction
    interval on these months. \\
    Same-month-last-year exists for only 38.4\% of rows, all from December 2012 on. &
    Use a time-based split with 2013 months in the test set; handle the missing lag explicitly. \\
    Weekends run 3.5\% and bank holidays 1.8\% above the household's average day. &
    Keep the remaining-weekend and remaining-holiday counts: weak, but known in advance and leak-free. \\
    \bottomrule
  \end{tabularx}
\end{table}

% =============================================================================
\section{Feature Engineering}
\label{sec:features}

\subsection{The \texorpdfstring{Day-$d$}{Day-d} Dataset}
\label{sec:dataset}

The modelling table has \textbf{one row per household $\times$ month $\times$ $d$}, for $d \in \{7, 14, 21\}$:
84,807 household-months $\times$ 3 = \textbf{254,421 rows}. One rule keeps it honest: \emph{a feature may use
only days $1..d$ of the target month, or earlier months}. Days after $d$ appear only in the targets.
Table~\ref{tab:features} lists the features.

\begin{table}[!ht]
  \centering
  \caption{Selected features, with description and justification. ``MTD'' = month to date (days $1..d$).}
  \label{tab:features}
  \footnotesize
  % Y{f} widths sum to 2 (two Y columns; the feature column is fixed width).
  \begin{tabularx}{\linewidth}{>{\RaggedRight\arraybackslash}p{3.7cm} Y{0.85} Y{1.15}}
    \tableheadrow
    \colhead{Feature} & \colhead{Description} & \colhead{Justification} \\
    \tableheadend
    \multicolumn{3}{l}{\textit{Month to date}} \\
    \code{kwh\_to\_date} & kWh used on days $1..d$ & Known part of the bill; $r = 0.98$ with the final bill at $d = 14$. \\
    \code{daily\_mean\_to\_date} & Mean daily kWh, days $1..d$ & Usage level independent of $d$ and month length. \\
    \code{last7\_mean\_kwh} & Mean daily kWh, days $d{-}6..d$ & Picks up a change in usage partway through the month; strongest single correlate of the remaining kWh at $d = 21$ ($r = 0.94$). Same as the MTD mean at $d = 7$. \\
    \code{run\_rate\_kwh} & \code{kwh\_to\_date}~$/\,d~\times$ days in month & Naive full-month projection; the benchmark the model must beat. \\
    \code{days\_elapsed}, \code{days\_remaining} & $d$, and days in month $- d$ & Scale the remaining kWh; account for 28 to 31-day months. \\
    \multicolumn{3}{l}{\textit{Earlier months}} \\
    \code{previous\_month\_kwh} & kWh in the previous calendar month (NULL after a gap) & Recent household level; most informative at small $d$. \\
    \code{rolling\_3month\_kwh} & Mean kWh of the three previous calendar months & Smoother level, less affected by one unusual month. \\
    \code{same\_month\_last\_year\_kwh} & kWh in the same month a year earlier & Captures the household's own seasonal pattern; available from December 2012. \\
    \multicolumn{3}{l}{\textit{Calendar (known in advance)}} \\
    \code{month\_sin}, \code{month\_cos} & Month on a circle: $\sin$ and $\cos$ of $2\pi(m-1)/12$ & Seasonality with December next to January; replaces month number, quarter and season dummies. \\
    \code{weekend\_days\_remaining} & Saturdays and Sundays after day $d$ & Weekend days use 3.5\% more than average. \\
    \code{holidays\_remaining} & Bank holidays after day $d$ & Holidays use 1.8\% more than average. \\
    \multicolumn{3}{l}{\textit{Weather and household}} \\
    \code{temp\_mean\_to\_date} & Mean daily temperature, days $1..d$ & Heating demand. Full-month temperature would leak days after $d$. \\
    \code{acorn\_comfortable}, \code{acorn\_adversity}, \code{acorn\_unknown} & ACORN group, one-hot, Affluent as the reference & Household level, useful when history is short. One group is dropped so the dummies do not sum to 1. \\
    \multicolumn{3}{l}{\textit{Targets and benchmarks (not inputs)}} \\
    \code{remaining\_kwh} & kWh on days $d{+}1$..end & Model target. \\
    \code{month\_kwh}, \code{final\_bill\_gbp} & Month total; total $\times$ \pounds 0.14228 & Final bill = \code{kwh\_to\_date} + predicted \code{remaining\_kwh}, times the tariff. \\
    \code{runrate\_bill\_gbp}, \code{last\_month\_bill\_gbp}, \code{same\_month\_ly\_bill\_gbp} & The three Phase~1 benchmarks as bills & Let Phase~4 score the benchmarks directly. \\
    \bottomrule
  \end{tabularx}
\end{table}

\subsection{Feature Availability}
\label{sec:availability}

Only the earlier-month features can be missing (Table~\ref{tab:availability}); all other columns are complete.
Missing lags are left as NULL rather than filled in, because filling them with, for example, the household mean
would use months after the target month. Phase~4 handles them either with models that accept missing values
(gradient-boosted trees) or with a missing-value indicator for linear models.

\begin{table}[!ht]
  \centering
  \caption{Availability of the earlier-month features (identical for every $d$; 84,807 household-months).}
  \label{tab:availability}
  \small
  \begin{tabularx}{\linewidth}{X r r r}
    \tableheadrow
    \colhead{Feature} & \colhead{Available} & \colhead{Missing} & \colhead{Available (\%)} \\
    \tableheadend
    \code{previous\_month\_kwh} & 80,139 & 4,668 & 94.50 \\
    \code{rolling\_3month\_kwh} & 71,002 & 13,805 & 83.72 \\
    \code{same\_month\_last\_year\_kwh} & 32,588 & 52,219 & 38.43 \\
    \bottomrule
  \end{tabularx}
\end{table}

\subsection{Relationships Among Features}
\label{sec:corr}

Figure~\ref{fig:corr} shows the correlations at $d = 14$. The four month-to-date features are nearly the same
signal ($r \geq 0.99$ with each other), so a linear model should keep one of them or use regularisation; tree
models are not harmed by the redundancy. For the remaining kWh, \code{previous\_month\_kwh} correlates at 0.910,
0.889 and 0.866 at $d = 7, 14, 21$, while \code{last7\_mean\_kwh} rises from 0.927 to 0.940: history matters
most early in the month. The month cosine and temperature to date both measure winter ($r = -0.89$). The weekend
and holiday counts are weak on their own ($|r| < 0.1$).

\begin{figure}[!ht]
  \centering
  \includegraphics[width=0.92\linewidth]{p3_feature_corr.png}
  \caption{Pearson correlations of the day-$d$ features and targets at $d = 14$ (pairwise on rows where both
  values exist).}
  \label{fig:corr}
\end{figure}

\subsection{Encoding and Scaling}
\label{sec:scaling}

ACORN group is one-hot encoded with Affluent, the largest group, as the reference; the month is encoded as a
sine/cosine pair; the tariff indicator from Phase~2 is dropped because every row is Std. Features sit on very
different ranges (\code{kwh\_to\_date} up to 3,741~kWh, the month encodings between $-1$ and 1), so scale-sensitive
models need scaling. Scaling is \textbf{deferred to Phase~4} and will be fitted on the training period only, so
the test months cannot influence it. Given the skew and outliers in Section~\ref{sec:outliers}, a robust scaler
(median and IQR) or a log transform of the kWh columns is worth testing.

\subsection{Changes from the First Pass and Validation}
\label{sec:validation}

The first pass predicted next month's bill from previous months. The changes made to realign it with the
Phase~1 day-$d$ task are:
\begin{itemize}[nosep]
  \item Std households only, filtered before any statistic or lag.
  \item New month-to-date features from the daily data, and a 12-month lag for the third benchmark.
  \item Month number, quarter and four season dummies replaced by one sine/cosine pair; one ACORN dummy dropped.
  \item Full-month temperature replaced by temperature to date, which does not leak future days.
  \item Pooled IQR outliers reinterpreted, and a within-household check added.
\end{itemize}

The notebook ends with automatic checks that stop it if any fails. All pass: no duplicate household $\times$
month $\times$ $d$ rows; the row count equals household-months $\times$ 3; \code{kwh\_to\_date} +
\code{remaining\_kwh} equals the month total for every row; the final bill equals the Phase~2 monthly bill; every
lag points to the correct calendar month. To test for leakage, every month-to-date feature was rebuilt from a copy
of the daily data cut off at day $d$, and all values match. One household-month was also recomputed by hand from
its daily rows at each $d$.

% =============================================================================
\section{Conclusion}
\label{sec:conclusion}

The Std households show strongly skewed, seasonal, temperature-driven consumption, with much larger differences
between households than within them. The daily data shows that by mid-month the naive run-rate projection already has a median error of 4.2\% on the month total, far better than the history-only benchmarks. The
modelling dataset (\code{features\_day\_d.parquet}, 254,421 rows) contains month-to-date, history, calendar,
weather and household features that use only information available on day $d$, plus the targets and all three
benchmarks. Phase~4 will model the remaining kWh, choose $d$, compare pooled and per-household models against the
run-rate projection, and fit scaling and prediction intervals on the training period only.

\printbibliography[title={References}]

\end{document}
